\section[召南]{召南}

\subsection[鵲巢]{鵲巢}

維鵲有巢，維鳩處之，之子于歸，百兩御之。

維鵲有巢，維鳩方之，之子于歸，百兩將之。

維鵲有巢，維鳩盈之，之子于歸，百兩成之。

\subsection[采蘩]{采蘩}

于以采蘩，于沼于沚，于以用之，公侯之事。

于以采蘩，于澗之中，于以用之，公侯之宫。

被之僮僮，夙夜在公，被之祁祁，薄言還歸。

\subsection[草蟲]{草蟲}

喓喓草蟲，趯趯阜螽，未見君子，憂心忡忡，亦既見止，亦既覯止，我心則降。

陟彼南山，言采其蕨，未見君子，憂心惙惙，亦既見止，亦既覯止，我心則説。

陟彼南山，言采其薇，未見君子，我心傷悲，亦既見止，亦既覯止，我心則夷。

\subsection[采蘋]{采蘋}

于以采蘋？南澗之濱，于以采藻？于彼行潦。

于以盛之？維筐及筥，于以湘之？維錡及釜。

于以奠之？宗室牖下，誰其尸之？有齊季女。

\subsection[甘棠]{甘棠}

蔽芾甘棠，勿翦勿伐，召伯所茇。

蔽芾甘棠，勿翦勿敗，召伯所憇。

蔽芾甘棠，勿翦勿拜，召伯所説。

\subsection[行露]{行露}

厭浥行露，豈不夙夜，謂行多露。

誰謂雀無角？何以穿我屋？誰謂女無家？何以速我獄？雖速我獄，室家不足。

誰謂鼠無牙？何以穿我墉？誰謂女無家？何以速我訟？雖速我訟，亦不女從。

\subsection[羔羊]{羔羊}

羔羊之皮，素絲五紽，退食自公，委蛇委蛇。

羔羊之革，素絲五緎，委蛇委蛇，自公退食。

羔羊之縫，素絲五緫，委蛇委蛇，退食自公。

\subsection[殷其靁]{殷其靁}

殷其靁，在南山之陽，何斯違斯，莫敢或遑，振振君子，歸哉歸哉。

殷其靁，在南山之側，何斯違斯，莫敢遑息，振振君子，歸哉歸哉。

殷其靁，在南山之下，何斯違斯，莫或遑處，振振君子，歸哉歸哉。

\subsection[摽有梅]{摽有梅}

摽有梅，其實七兮，求我庶士，迨其吉兮。

摽有梅，其實三兮，求我庶士，迨其今兮。

摽有梅，頃筐塈之，求我庶士，迨其謂之。

\subsection[小星]{小星}

嘒彼小星，三五在東，肅肅宵征，夙夜在公，寔命不同。

嘒彼小星，維參與昴，肅肅宵征，抱衾與裯，寔命不猶。

\subsection[江有汜]{江有汜}

江有汜，之子歸，不我以，不我以，其後也悔。

江有渚，之子歸，不我與，不我與，其後也處。

江有沱，之子歸，不我過，不我過，其嘯也歌。

\subsection[野有死麕]{野有死麕}

野有死麕，白茅包之，有女懷春，吉士誘之。

林有樸樕，野有死鹿，白茅純束，有女如玉。

舒而脱脱兮，無感我帨兮，無使尨也吠。

\subsection[何彼襛矣]{何彼襛矣}

何彼穠矣，唐棣之華，曷不肅雝，王姬之車。

何彼穠矣，華如桃李，平王之孫，齊侯之子。

其釣維何，維絲伊緡，齊侯之子，平王之孫。

\subsection[騶虞]{騶虞}

彼茁者葭，壹發五豝，于嗟乎騶虞。

彼茁者蓬，壹發五豵，于嗟乎騶虞。
